\documentclass{article}
\usepackage{amsmath}
\usepackage{amssymb}
\usepackage{amsthm}
\usepackage{mathtools}
\usepackage{geometry}
\usepackage{hyperref}

\newcommand{\rk}{\operatorname{rank}}
\newcommand{\GL}{\operatorname{GL}}
\newcommand{\R}{\mathbb{R}}
\newcommand{\Mat}{\mathrm{Mat}}

\theoremstyle{definition}
\newtheorem{theorem}{Theorem}[section]
\newtheorem{proposition}[theorem]{Proposition}
\newtheorem{lemma}[theorem]{Lemma}
\newtheorem{corollary}[theorem]{Corollary}
\newtheorem{definition}[theorem]{Definition}
\newtheorem{remark}[theorem]{Remark}
\newtheorem{conjecture}[theorem]{Conjecture}

\title{Task 5: Proof of Conjecture A}
\date{}

\begin{document}
\maketitle

\section{Goal}

Deduce Conjecture A from the normalization theorem proved in Task 4.

\begin{conjecture}[Conjecture A]
For every point $W\in \Gamma_{d,S}^r$, there exists an element
\[
g\in G_d^{W_0}
\]
such that
\[
\mu(g\cdot W)=P_SW^\ast.
\]
Equivalently, every $G_d^{W_0}$-orbit in $\Gamma_{d,S}^r$ meets the critical locus $\mathcal C_{d,S}^r$.
\end{conjecture}

\section{Required proof structure}

\begin{enumerate}
\item start from $W\in \Gamma_{d,S}^r$;
\item apply the normalization theorem to $\mu(W)$;
\item choose an element of $G_d^{W_0}$ realizing the endpoint normalization;
\item prove that the transformed tuple remains in $\Gamma_{d,S}^r$;
\item conclude that the transformed tuple lies in $\mathcal C_{d,S}^r$.
\end{enumerate}

\section{Deliverable}

This file is complete once Conjecture A is fully proved.

\end{document}
